\documentclass[11pt]{article}
\usepackage[a4paper,margin=1in]{geometry}
\usepackage{amsmath,amssymb,amsthm,booktabs}
\usepackage[T1]{fontenc}
\usepackage{lmodern}
\usepackage[hidelinks]{hyperref}
\setlength{\emergencystretch}{3em}

\theoremstyle{plain}
\newtheorem{theorem}{Theorem}
\newtheorem{lemma}[theorem]{Lemma}
\newtheorem{proposition}[theorem]{Proposition}
\theoremstyle{remark}
\newtheorem{remark}[theorem]{Remark}

\title{A Local-Obstruction Counterexample to Conjecture 00000000008}
\author{%
  Feiyu\thanks{Independent researcher.}
  \and
  Claude (Opus 5)\thanks{Anthropic.}%
}
\date{2026-09-12}

\begin{document}
\maketitle

\begin{abstract}
Conjecture 00000000008 asserts that the prime-value count of a square-free
integer-coefficient quadratic $f$ admits a lower bound $c_f\,x/\log x$ with
$c_f>0$. As stated it omits the condition that no prime divides every value of
$f$, and it is false: $f(n)=n^2+n+2$ is irreducible over $\mathbb Q$, hence
square-free, yet every value is even and its only prime value is $f(0)=2$. The
refutation is formalized in Lean~4 against Mathlib.
\end{abstract}

\section{Statement}

Conjecture 00000000008 concerns the number of prime values of a square-free
integer-coefficient quadratic polynomial $f$, and asserts that there is a
constant $c_f>0$ with
\begin{equation}\label{eq:bound}
  \Pi_f(x)\;\ge\;c_f\,\frac{x}{\log x}
\end{equation}
for all large $x$, where $\Pi_f(x)=\#\{n\le x:f(n)\text{ is prime}\}$.

The conditions imposed are that $f$ be quadratic, have integer coefficients,
and be square-free. What is \emph{not} imposed --- and what the classical
Bunyakovsky and Hardy--Littlewood formulations do impose --- is that no prime
divide every value of $f$. Without it the assertion is false.

\section{The counterexample}

\begin{proposition}\label{prop:hyp}
Let
\[
  f(n)=n^2+n+2 .
\]
Then $f$ is the value function of the integer quadratic
\[
  F=X^2+X+2,
\]
which
\begin{enumerate}
  \item has integer coefficients $1,1,2$ and degree $2$;
  \item has discriminant $1^2-4\cdot1\cdot2=-7\ne0$, so no repeated root;
  \item is irreducible over $\mathbb Q$, hence square-free;
  \item has positive leading coefficient.
\end{enumerate}
\end{proposition}

\begin{proof}
Only irreducibility needs an argument. For every $x\in\mathbb Q$,
\[
  4\left(x^2+x+2\right)=(2x+1)^2+7\;\ge\;7\;>\;0,
\]
so $F$ has no rational root. A polynomial of degree two over a field with no
root in that field is irreducible; and an irreducible element is square-free.
\end{proof}

So $f$ meets every condition the conjecture states, and the classical positivity
and non-degeneracy conditions besides.

\begin{theorem}\label{thm:main}
$f(0)=2$ is the only prime value of $f$. Consequently $\Pi_f(x)=1$ for every
$x\ge0$, and no constant $c_f>0$ satisfies \eqref{eq:bound}.
\end{theorem}

\begin{proof}
Write $f(n)=n(n+1)+2$. One of $n$, $n+1$ is even, so $n(n+1)$ is even and
$f(n)$ is even for every $n$. For $n\ge1$ we have $n\le n^2$, hence
\[
  f(n)=n^2+n+2\;\ge\;4 .
\]
So for $n\ge1$ the value $f(n)$ is an even integer greater than $2$, therefore
composite. At $n=0$ we get $f(0)=2$, which is prime. Hence
$\Pi_f(x)=1$ for every $x\ge0$.

Suppose \eqref{eq:bound} held with some $c_f>0$ for all $x\ge N$. In the
equivalent division-free form this says
\[
  c_f\,x\;\le\;\Pi_f(x)\log x\;=\;\log x
  \qquad (x\ge N).
\]
But $\log x=o(x)$, so $\log x\le\tfrac{c_f}2x$ for all sufficiently large $x$.
Choosing $x$ at least $N$, at least $1$, and large enough for that estimate
gives $c_f x\le\tfrac{c_f}2x$, hence $c_f x\le 0$, contradicting $c_f>0$ and
$x\ge1$.
\end{proof}

\section{What the counterexample shows}

The obstruction is local, at the prime $2$:
\[
  f(n)\equiv0\pmod2\qquad\text{for every integer }n .
\]
It is worth being precise about which hypothesis fails, because $f$ satisfies
all the others. It is quadratic, has integer coefficients, has nonzero
discriminant, is irreducible over $\mathbb Q$ --- so it is square-free under
either reading of that word --- and has positive leading coefficient. The sole
failure is the fixed prime divisor $2$.

A repaired conjecture must therefore add the Bunyakovsky condition: the values
$f(n)$, $n\in\mathbb N$, have greatest common divisor $1$. Under that condition
the assertion becomes the quadratic case of the Bunyakovsky conjecture, which
is open --- it is not known that any irreducible quadratic takes infinitely
many prime values.

\begin{remark}
Squarefreeness is not what is doing the work here, and cannot be made to do it:
one can fail the conjecture with a polynomial that is irreducible, which is
strictly stronger than square-free. Nor is reducibility needed; $n(n+1)$ would
also give a counterexample, but it is reducible and so a referee could object
that ``square-free'' was meant to exclude it. The polynomial above forecloses
that objection.
\end{remark}

\section{Formal verification}

The accompanying Lean~4 project formalizes the refutation against Mathlib. The
conditions the conjecture places on $f$ are established \emph{about $f$}, not
as side computations on literals: \texttt{F\_eval} identifies $f$ with the
value function of the polynomial \texttt{F}, and every subsequent fact about
\texttt{F} therefore speaks about $f$.

\begin{itemize}
  \item \texttt{F\_eval} --- $\texttt{F.eval}\,n=f(n)$ for every $n$, where
    $\texttt{F}=X^2+X+2$ over $\mathbb Q$.
  \item \texttt{F\_natDegree}, \texttt{F\_coeff\_zero/one/two},
    \texttt{F\_discriminant} --- degree $2$, coefficients $2,1,1$, and
    discriminant $-7$.
  \item \texttt{F\_irreducible}, \texttt{F\_squarefree} --- irreducibility via
    Mathlib's \texttt{irreducible\_of\_degree\_le\_three\_of\_not\_isRoot},
    and \texttt{Squarefree F} is Mathlib's \texttt{Squarefree}, not a
    re-definition.
  \item \texttt{f\_even}, \texttt{f\_ge\_four}, \texttt{f\_not\_prime} --- the
    elementary facts, with primality Mathlib's \texttt{Nat.Prime}.
  \item \texttt{f\_prime\_iff} --- $\texttt{Nat.Prime}(f\,n)\iff n=0$, the
    complete description of the prime values.
  \item \texttt{primeValueCount\_eq\_one} --- $\Pi_f(x)=1$ for every $x$, the
    count taken over all $n\le x$ including $n=0$.
  \item \texttt{conjecture\_00000000008\_false} ---
    $\lnot\,\texttt{LowerBoundHolds}$, where \texttt{LowerBoundHolds} states
    \eqref{eq:bound} in its division-free form. The proof uses Mathlib's
    \texttt{Real.isLittleO\_log\_id\_atTop}.
\end{itemize}

Because the conjectured lower bound is itself stated in the formalization, the
final theorem is a formalization of its negation rather than of facts that
merely entail it.

The project builds with \texttt{lake build} on
\texttt{leanprover/lean4:v4.33.1} with Mathlib \texttt{v4.33.1}. It contains no
\texttt{sorry}, no \texttt{admit} and no \texttt{native\_decide}, and
introduces no axioms of its own:
\texttt{\#print axioms conjecture\_00000000008\_false} reports exactly the
three standard axioms \texttt{propext}, \texttt{Classical.choice} and
\texttt{Quot.sound}.

\end{document}
